\section{Introduction}\label{s:intro}
Can a fixed local noisy quantum medium be an excellent memory at one
noise rate and destroy every encoding at another, with no algorithm to
distinguish its computably named rate windows? We construct such a
medium, separately for two physical noise models. Nonhalting windows
store a reference-entangled qubit for an exponential number of rounds;
halting windows destroy every encoding after a time independent of size.
The two media occupy six dimensions and use seven elementary gates.
The three-dimensional construction retains a sharp linear storage scale.
A further whole-site medium gives stretched-exponential storage on the
undecidable windows in four dimensions.
Lean 4 with Mathlib verifies the Bell-sharing bound
(Lemma~\ref{f:bell}) and conditional window certification
(Theorem~\ref{lean:conditional}); Appendix~\ref{app:lean} gives the scope.

Fix a universal machine \(U\), and write \(U(n)\uparrow\) or
\(U(n)\downarrow\) for nonhalting or halting on input \(n\). The precise
model and fidelity convention are in Section~\ref{s:model}. Unless specified as elementary-layer time, time counts repetitions of
one fixed finite circuit, called a round.

The original three-dimensional construction has a sharp lifetime scale.
\begin{maintheorem}\label{thm:main}
For each of whole-site and independent single-qubit depolarization after
every layer, there is a separately constructed, computably specified,
translation-invariant, time-periodic finite-range circuit
on \(\Z^3\), with gate inventory
\[
 \{X,H,\mathrm{CNOT},\mathrm{Toffoli},R_0,D_Z,I\}
\]
and computable disjoint rational windows \(I_n\): nonhalting gives optimal storage time
\(\Theta(L)\); halting gives a time bound uniform in \(L\) and the window;
certification of its explicit linear-time storage circuits is
\(\Pi^0_1\)-complete.
\end{maintheorem}

Exponential storage is possible on the undecidable windows themselves.
\begin{maintheorem}\label{thm:exponential}
For each of whole-site and independent single-qubit depolarization after
every elementary layer, a separately compiled, computably specified,
translation-invariant, time-periodic finite-range circuit on \(\Z^6\)
with the same seven gates has computable disjoint rational windows
\(I_n\). For one rational \(c_{\mathrm e}>0\), independent of \(n\)
and the same for both media:
if \(U(n)\uparrow\), one encoder keeps entanglement fidelity at least
\(15/16\) at every rate in \(I_n\) and every time
\(t\le2^{\lfloor c_{\mathrm e}L\rfloor}\), for all \(L\) above a
computable threshold \(L_{\mathrm e}(n)\); if \(U(n)\downarrow\), every
encoding is destroyed after a time uniform in \(L\) and the window.
\end{maintheorem}

The encoder and the time-dependent decoder are total computable and do
not depend on the rate or the halting answer. The detailed statement is
Theorem~\ref{v:detailed}. Proposition~\ref{v:certification} adds a
coarse \(\exp(O_n(L^4))\) lifetime upper bound and exact
\(\Pi^0_1\)-complete certification for these exponential media.
The matching linear upper bound applies to Theorem A.

Beyond-polynomial storage on the undecidable windows requires only four dimensions under whole-site noise.
\setcounter{maintheorem}{3}
\begin{maintheorem}\label{thm:four-dimensional}
For whole-site depolarization after every elementary layer, a computably
specified, translation-invariant, time-periodic finite-range circuit on
\(\Z^4\) with the same seven gates has computable disjoint rational
windows \(I_n\). For rationals \(c_4,\alpha_4>0\), independent of \(n\):
if \(U(n)\uparrow\), one encoder keeps entanglement fidelity at least
\(15/16\) at every rate in \(I_n\) and every elementary-layer time
\(t\le\exp(c_4L^{\alpha_4})\), for all \(L\) above a computable threshold
\(L_4(n)\); if \(U(n)\downarrow\), every encoding is destroyed after a
time uniform in \(L\) and the window.
\end{maintheorem}
Here \(L\) is the torus side length. We call it Theorem D because
Appendix C already contains Proposition C.1, a different four-dimensional
result. Its encoder and decoders are total computable and independent
of the rate and the halting answer.
The detailed scope is Theorem~\ref{f4:detailed}. Its exponent is
astronomically small, below \(2\cdot10^{-4}\); no lifetime upper bound
is proved for this medium.



For Theorem A, storage time is the last time at which some encoder and decoder reach
entanglement fidelity \(7/8\). If \(U(n)\) never halts, then for every
side \(L\ge L_0(n)\), with \(L_0\) computable, the explicit construction
keeps fidelity at least \(15/16\) through time \(\lfloor cL\rfloor\).
The rational \(c>0\) does not depend on \(n\), and neither does the
time \(201L+C_0\) of the matching upper bound. On the halting side, the optimal fidelity
is at most \(1/2+\eta\) after a time uniform in both volume and the
chosen window, for every \(\eta>0\). The detailed theorem,
Theorem~\ref{thm:detailed}, also gives a whole-channel upper bound.

Finite systems forget in both cases. Independent full-support noise has
a positive chance of replacing every site in one layer; this eventually
removes the input. The distinction is the size dependence of useful
storage time. For Theorem A, linear storage is optimal even
though the constants in the two bounds differ enormously. All three theorems are existence results: their local dimensions,
interaction ranges and noise scales are astronomical. Each site carries a fixed number of
qubits, and one layer may apply many gates of at most three qubits at
the same site, on disjoint qubits; the noise then acts once per site
(\(\mathrm W\)) or once per qubit (\(\mathrm Q\)) after every layer.
Encoding and final decoding are global,
noiseless, computable operations; the repeated correction inside the
medium is local.

\begin{figure}[!ht]
\centering
\begin{tikzpicture}[x=1cm,y=1cm,>=Stealth,font=\small]
\draw[->] (0,0)--(5.8,0) node[right]{\(\eps\)};
\foreach \x/\w/\lab in {0.45/0.18/n+2,1.6/0.28/n+1,2.9/0.45/n,4.6/0.7/1}{
 \fill[blue!18] (\x,0.1) rectangle (\x+\w,0.7);
 \draw[blue!60!black] (\x,0.1) rectangle (\x+\w,0.7);
 \node[above] at (\x+\w/2,0.7) {\(I_{\lab}\)};}
\node[align=center] at (2.8,-0.55){Computably named noise windows};
\begin{scope}[shift={(7,0)}]
 \draw[->] (0,0)--(4.8,0) node[right]{\(L\)};
 \draw[->] (0,0)--(0,2.5) node[above]{storage time};
 \draw[thick,blue!65!black] (0.35,0.25) .. controls (2.5,0.35) and (3.3,0.7) .. (3.9,2.2);
 \node[above] at (2.4,1.65) {\(U(n)\uparrow:\ \ge2^{\lfloor c_{\mathrm e}L\rfloor}\)};
 \draw[thick,red!65!black] (0.35,0.65)--(4.1,0.65);
 \node[above] at (2.9,0.65) {\(U(n)\downarrow:\ O_n(1)\)};
\end{scope}
\node[draw,rounded corners,align=center] (a) at (1.2,-1.7){static noisy chart\\and finite counter};
\node[draw,rounded corners,align=center] (b) at (5.4,-1.7){local toric-code\\correction};
\node[draw,rounded corners,align=center] (c) at (9.7,-1.7){final controlled\\erasure};
\draw[->] (a)--(b); \draw[->] (b)--(c);
\draw[->] (a.south) to[bend right=18] (c.south);
\end{tikzpicture}
\caption{Theorem B: exponential guaranteed storage or size-independent
destruction on computable noise windows. Locations, widths and constants
are schematic; the blue curve is a lower bound. Theorem A separately
has the sharp linear scale. Both constructions end each round with a
complete permission-controlled cut.}
\label{fig:story}
\end{figure}

\noindent\textbf{How the medium works.}
For Theorem B, the classical computation is a static chart written once
into a square of side \(\lfloor L/4\rfloor\) and separated from its
periodic copies by vacuum, so it never has to wrap around the torus. A
binary counter along the square's diagonal grants permission to store
for \(2^{\lfloor L/4\rfloor}-1\) rounds. Four further coordinates carry a
closed four-dimensional toric code with a local sweep decoder. Controller
errors are correlated, so we bound them by counting complete witnesses
of each wrong erasure, keeping enough coordinates to tell distinct noise
events apart on the torus. On a halting window every region of the chart
contains a local violation or a wrong rate test, or lies in a square small
enough that its counter expires; permission then dies, and with it every
encoding.

For Theorem D, the same static cap occupies two coordinates and is
protected in two transverse directions. A separate protected classical calendar schedules
a concatenated quantum code along a chain. Folding its cycle into a
segment uses two semantic qubits per site and keeps the interaction range fixed.
The code hierarchy reaches a fixed fraction of the system size. Short
error explanations have consistent source names on the torus; longer
ones have an exponentially small probability. This yields a
stretched-exponential storage guarantee in four dimensions.

For Theorem A, each site runs a noisy classical computation beside a quantum chain.
A second, transversely protected classical field supplies the calendar
for concatenated quantum error correction; higher clocks and programs
remain classical. The computation never reads the noise rate.
Instead, local counters compare how often fresh noise marks appear along
rows of lengths \(2\cdot4^r-1\). On the window \(I_n\) this selects the
depth at which the computation simulates \(U\) on input \(n\). Every
round ends with a controlled erasure of the quantum data at each site
where the computation is not live. For each nonhalting \(n\) and all
\(L\ge L_0(n)\), the computation stays live near the code long enough
for local correction to preserve the stored qubit through
\(\lfloor cL\rfloor\). If \(U(n)\) halts, every chart contains a wrong
gate, and the resulting bounds on paths of quantum dependence destroy
every encoding after a time uniform in system size. Even without halting, the computation needs an
aperiodic chart, which no side-\(L\) torus can carry. By \(200L\) rounds,
continued liveness requires exceptional noise events. Their path bounds
cut the remaining quantum dependence by \(201L+C_0\), outside an event
of exponentially small probability in \(L\). This gives the linear
upper bound and makes the storage scale optimal.

\subsection*{Context and credit}
Undecidable properties of quantum many-body systems include the spectral
gap \cite{CPGW2015} and ground-state phase diagrams. Bausch, Cubitt and
Watson \cite[Theorem~2.1]{BCW2021} construct one continuously
parameterized Hamiltonian with machine-input intervals: nonhalting
gives a gapped product phase on a fixed interval, while halting gives
a critical gapless phase on an interval whose width is specified
existentially. Purcell, Li and Cubitt
\cite{PLC2024} place a single phase boundary at an uncomputable point. Quantum-channel fidelity reachability was already shown
undecidable by Wolf, Cubitt and P\'erez-Garc\'ia \cite{WCPG2011}.
Table~\ref{tab:landscape} locates the operational question studied here.
Our parameter is the rate of independent physical noise following a
fixed circuit, and our negative branch quantifies over encodings into
all its registers. Here the windows are computable on both branches.

\begin{table}[t]
\caption{Closely related results and the observable studied here.}
\label{tab:landscape}
\small\centering
\begin{tabular}{@{}>{\raggedright\arraybackslash}p{0.29\textwidth}>{\raggedright\arraybackslash}p{0.29\textwidth}>{\raggedright\arraybackslash}p{0.34\textwidth}@{}}
\toprule
Work & Setting & Observable\\
\midrule
Cubitt--P\'erez-Garc\'ia--Wolf; Bausch--Cubitt--Watson; Purcell--Li--Cubitt
 & Local Hamiltonians; interaction parameter & Gap and ground-state phase\\[3pt]
Wolf--Cubitt--P\'erez-Garc\'ia & Finite channel collections & Fidelity reached by some channel word\\[3pt]
D\"unnweber--Styliaris--Trivedi; Drechsler--Oreg--Aharonov; Balasubramanian--Davydova--Lake
 & Autonomous noisy quantum memories & Positive storage guarantees\\[3pt]
G\'acs--Reif; G\'acs; Aharonov--Ben-Or; Harrington
 & Hierarchical reliable computation and quantum correction & Protected schedules and encoded gates\\[3pt]
Douglas--Mghirbi & Noisy classical local rules & Ergodicity and computable noise-window selection\\[3pt]
This paper, Theorem A & Fixed seven-gate circuits on \(\Z^3\); physical noise rate & Sharp linear storage, bounded forgetting, exact certification\\[3pt]
This paper, Theorem B & Separate seven-gate W/Q circuits on \(\Z^6\); physical noise rate & Exponential storage lower bound versus size-uniform forgetting on effective windows\\[3pt]
This paper, Theorem D & Seven-gate W circuit on \(\Z^4\); physical noise rate & Stretched-exponential storage lower bound versus size-uniform forgetting on effective windows\\
\bottomrule
\end{tabular}
\end{table}

Autonomous quantum memory under noisy local dynamics predates this
construction, notably the fixed-rule two-dimensional quantum automata
of D\"unnweber, Styliaris and Trivedi \cite{DST2026}, which already
combine G\'acs-style hierarchical classical control with a
measurement-free concatenated quantum code, and the local
toric-code memories of Balasubramanian, Davydova and Lake
\cite{BDL2026}. Both give lifetimes that grow without bound with system
size, in two or three dimensions. The construction of Theorem~A also
combines hierarchical reliable classical computation
\cite{Gacs2001,GacsReif1988} with concatenated quantum fault tolerance
\cite{ABO1999,ABO2008}. The additional conclusions here are an
undecidable noise-window structure with destruction for every encoding;
for Theorem~A, a sharp linear storage scale and exact certification;
for Theorem~B, exponential storage on the undecidable windows
themselves; for Theorem~D, a stretched-exponential guarantee there in
four dimensions under whole-site noise.
Harrington's hierarchical local decoder
\cite{Harrington2004} is an early quantum-memory precedent. The exponential constructions use the toric-code and sweep lineage of
Dennis, Kitaev, Landahl and Preskill \cite{DKLP2002}, Ahn
\cite{Ahn2004}, Kubica and Preskill \cite{KubicaPreskill2019}, and
Vasmer, Browne and Kubica \cite{VBK2021}. Particularly close comparisons are
the classical-to-quantum memory correspondence of Drechsler, Oreg and
Aharonov \cite{DOA2026} and the cubical sweep fault-tolerance theorem of
Ponarovsky \cite{Ponarovsky2026}. We prove the source-dependent estimate
needed for this medium rather than invoking a threshold for independent
faults: the controller induces correlated erasures.

The classical computational controller builds on our preprints on
fixed-noise ergodicity, undecidable phase diagrams, and an uncomputable
noise diagram \cite{DMFixed2026,DMPhase2026,DMNoise2026}. Those works
supply the finite computational charts and classical probability
estimates. Version 1.0.0 of this paper \cite{DMQuantum2026} already proved the
linear quantum bridge and a designed exponential window. Version 2.0.0 \cite{DMQuantumV200} already supplies the static cap,
canonical initializer, correlated-source method and six-dimensional
Theorem B, together with Section~\ref{sec:quantum}'s real code gadgets
and Section~\ref{sec:forgetting}'s full-channel sharing argument.
Theorem D combines those published ingredients, including the
private-port dispatcher of Section~\ref{d:dispatch} and the fold of
virtual cycles into physical sheets with sheet-typed fault labels from
Appendices~\ref{sec:fold-controls}--\ref{sec:folded}, with whole-state
periodic calendar resets, exclusive port selection and a torus source
analysis that controls aliases of the same physical fault. Its new
conclusion is beyond-polynomial storage on the undecidable windows in
four dimensions for whole-site noise.

Appendix~\ref{sec:coexistence} retains a separate four-dimensional
construction. Proposition~\ref{p:coexistence} gives one independent-qubit
medium with the same seven gates, an uncomputable storage diagram, and a
reserved, designed and decidable window storing a qubit for
\(2^{\lfloor c_{\exp}L\rfloor}\); its universal nonhalting windows have
linear storage. Version 1.0.0 left exponential storage on the undecidable
windows themselves open; Theorem B provides it, in six dimensions and for
both noise models. Its proof uses the deterministic finite-torus
machinery of Lemma~\ref{x:exponential}.

\paragraph{Changes in 2.0.0.}
Theorem B adds exponential storage on the undecidable windows for
separate six-dimensional media in both noise models. Theorem A and
the four-dimensional designed-window result are retained. The new
proof, a coarse lifetime upper bound and exponential-horizon
certification are in Section~\ref{sec:exponential}; the discussion
states the remaining dimensional and matching-lifetime questions.

\paragraph{Changes in 2.1.0.}
Theorem D lowers the dimension of beyond-polynomial storage on the
undecidable windows from six to four for whole-site noise.
Section~\ref{sec:four-dimensional} gives the construction and proof;
Appendices~\ref{app:f4-calendar}--\ref{app:f4-sources} prove the added
calendar, folded-dispatch and torus-source lemmas. Theorems A and B and
Appendices A--C retain their previous mathematical statements.
Appendix~\ref{app:lean} adds Lean verification of Lemma~\ref{f:bell}
and the printed conditional abstraction of Theorem~\ref{a:complete}.

\subsection*{Reading guide}
Section~\ref{s:model} states the model and full quantifiers.
Sections~\ref{sec:controller} and~\ref{sec:compilation} give the classical
selector and exact physical implementation. Section~\ref{sec:quantum}
proves three-dimensional storage, including the gadget support bound,
source counting and decoding at every requested time.
Sections~\ref{sec:forgetting} and~\ref{sec:ceiling} prove the two
all-encoding bounds; Section~\ref{a:certification} proves certification.
Section~\ref{sec:exponential} proves Theorem B, including its static
initializer, physical-name density bound, source law, quantum recovery
and full-algebra negative branch. Computability readers can proceed
from Theorem A to certification, or from Theorem B to its static
alternative and final channel argument.
Quantum-information readers can take the explicitly stated classical
chart interface as an input. The folded-code and exponential-coexistence
proofs are in Appendices~\ref{sec:fold-controls}--\ref{sec:coexistence}.

Section~\ref{sec:four-dimensional} can be read after the code and
normalizer lemmas of Section~\ref{sec:quantum} and the static cap of
Section~\ref{sec:exponential}. Its dependency tables specify the
constant order; its appendices contain the new source and fold proofs.
